% =====================================================================
%  第 II 部：仕組みから手法へ
%  文の元：paper/sections/5-part2-mechanism-to-method.tex（図 method，ladder）．基本版と拡張版の書き方は paper/CLAIM.md 第 5 版
%  図の値：ladder_levels.csv，ladder_gains.csv，capture.csv．弱いラベルの表の数は説明用の値（実験の値ではない）
% =====================================================================
\setband{第 II 部　仕組みから手法へ}

\begin{frame}{手法の全体：出発点を作り，先生の弱いラベルと $K$ 本で生徒を学ぶ}
  \rl{手法}
  \vspace{-0.3em}
  \centering
  \begin{tikzpicture}[font=\scriptsize, x=1cm, y=1cm]
    \fill[cSrc!45, rounded corners=3pt] (-0.2,2.35) rectangle (15.2,4.7);
    \fill[cTgt!45, rounded corners=3pt] (-0.2,-0.2) rectangle (15.2,2.15);
    \node[anchor=north west, font=\scriptsize\bfseries] at (-0.1,4.65) {ソース（ターゲットのコーパス以外の群）};
    \node[anchor=north west, font=\scriptsize\bfseries] at (-0.1,2.1) {ターゲット};
    \node[srcbox, text width=2.0cm] (s0) at (1.2,3.45) {ラベル付きの発話};
    \node[clf, text width=1.9cm] (sc) at (4.0,3.45) {中心化\\（ソースごと）};
    \node[draw=cStart, fill=white, rounded corners=2pt, align=center, text width=3.2cm, line width=0.8pt] (st) at (7.6,3.95) {\textbf{出発点} $(W_0, b_0)$\\生徒のモダリティ};
    \node[draw=cExt, fill=white, rounded corners=2pt, align=center, text width=3.2cm, line width=0.8pt] (te) at (7.6,2.95) {\textbf{先生}：融合の出発点など\\（同じ作り方，ソースだけ）};
    \draw[arr] (s0) -- (sc);
    \draw[arr] (sc.east) -- (st.west);
    \draw[arr] (sc.east) -- (te.west);
    \node[tgtbox, text width=2.0cm] (tp) at (1.2,0.95) {プール\\（テスト以外）};
    \node[clf, text width=1.9cm] (tc) at (4.0,0.95) {中心化\\（プールの平均）};
    \node[tgtbox, text width=2.4cm] (kl) at (7.6,1.55) {$K$ 本のラベル（$4K$ 本）};
    \node[tgtbox, text width=2.4cm] (ul) at (7.6,0.4) {ラベル無しの発話};
    \draw[arr] (tp) -- (tc);
    \draw[arr] (tc.east) -- (kl.west);
    \draw[arr] (tc.east) -- (ul.west);
    \node[draw=cExt, fill=cExt!8, rounded corners=2pt, align=center, text width=2.6cm] (wl) at (11.0,0.4) {\textbf{弱いラベル}\\感情ごとの合計をそろえ，\\全体を $m$ にする};
    \draw[arr, cExt] (ul) -- node[below, font=\tiny, text=cExt] {先生の予測} (wl);
    \draw[arr, cExt] (te.east) .. controls (10.4,2.95) and (11.0,2.2) .. (wl.north);
    \node[draw=cOurs, fill=white, rounded corners=2pt, align=center, text width=2.2cm, line width=0.9pt] (stu) at (13.9,1.55) {\textbf{生徒}（線形）\\$K$ 本と弱いラベル\\で学ぶ};
    \draw[arr, cOurs] (kl.east) -- (stu.west);
    \draw[arr, cExt] (wl.east) -- (stu.south west);
    \draw[cStart, line width=0.9pt, decorate, decoration={coil, aspect=0.4, segment length=3pt, amplitude=2.2pt}] (st.east) -- (stu.north);
    \node[font=\tiny, text=cStart, anchor=south west] at (11.6,3.3) {出発点に縛る};
  \end{tikzpicture}\par
  \vspace{0.2em}\raggedright
  \begin{itemize}\scriptsize\setlength{\itemsep}{2pt}
        \item 重みと切片は別の強さで出発点に縛る．5 つの値は仮のターゲットで 1 度，先生と $m$ は $K$ ごとに選ぶ
      \end{itemize}
  \hosoku{s:p2:obj}{生徒の目的関数}
\end{frame}

\begin{frame}{出発点の工夫ごとに UAR が上がる：仮説転移，さらに 5 つの値を選ぶ}
  \rl{結果：出発点}
  \vspace{-0.2em}
  \begin{columns}[T,onlytextwidth]
    \begin{column}{0.5\textwidth}
      \centering
      \begin{tikzpicture}
        \begin{axis}[width=7.0cm, height=4.8cm, ybar, bar width=5.5pt, ymin=30, ymax=62, xmin=0.5, xmax=4.5,
            xtick={1,...,4}, xticklabels={1,5,10,25}, xlabel={感情あたりのラベル $K$}, ylabel={UAR（\%）},
            ymajorgrids, grid style={mMuted!25},
            legend style={font=\tiny, draw=none, at={(0.5,1.02)}, anchor=south, legend columns=2, /tikz/every even column/.append style={column sep=0.6em}}]
          \addplot[fill=cTgt!90!black, draw=none] table[x=x, y=target] {ladder_levels.csv}; \addlegendentry{ターゲットのみ}
          \addplot[fill=mMuted!60, draw=none] table[x=x, y=pooled] {ladder_levels.csv}; \addlegendentry{合併（1 ソースの平均）}
          \addplot[fill=cStart, draw=none] table[x=x, y=p4] {ladder_levels.csv}; \addlegendentry{全ソースからの仮説転移}
          \addplot[fill=cOurs, draw=none] table[x=x, y=e39] {ladder_levels.csv}; \addlegendentry{5 つの値を選んだ出発点}
        \end{axis}
      \end{tikzpicture}\par
      {\raggedright\tiny A（WavLM）の生徒．ターゲット 7 つ，反復 30 回の平均\par}
    \end{column}
    \begin{column}{0.48\textwidth}
      \begin{itemize}\footnotesize\setlength{\itemsep}{10pt}
        \item 生徒は中心化した特徴で学ぶ（判断基準の損の大半が消える）
        \item 全ソースからの仮説転移は，従来法の最良を音声の $K \le 10$ で \NUM{p2-e31-p4-vs-best-audio-lo}〜\NUM{p2-e31-p4-vs-best-audio-hi} \pt 上回る
        \item 5 つの値を選ぶと，$K=5$ でさらに \NUM{p2-e39-vs-p4-k5-lo}〜\NUM{p2-e39-vs-p4-k5-hi} \pt 上がる
      \end{itemize}
    \end{column}
  \end{columns}
  \hosoku{s:p2:start}{出発点の作り方の式と全部の値}
\end{frame}

\begin{frame}{$K$ 本のラベルで直せる向きは，その $4K$ 本が張る空間の中だけである（命題 1）}
  \rl{考え方：向きの情報が得られないところ}
  \vspace{0.1em}
  \begin{columns}[T,onlytextwidth]
    \begin{column}{0.44\textwidth}
      \centering
      \begin{tikzpicture}[font=\scriptsize, scale=1.0]
        \fill[cOurs!10] (-0.6,-0.45) -- (4.0,-0.45) -- (4.9,1.1) -- (0.3,1.1) -- cycle;
        \node[text=cOurs, anchor=west, font=\tiny] at (2.6,-0.25) {$4K$ 本の点が張る空間};
        \fill[black] (0,0) circle (1.3pt) node[below left, font=\tiny] {$W_0$};
        \draw[->, line width=1.2pt, cBest] (0,0) -- (3.2,2.6) node[above, font=\tiny, text=cBest] {$E_0 = W^{*} - W_0$};
        \draw[->, line width=1.2pt, cOurs] (0,0) -- (3.2,0.55);
        \draw[densely dashed, mMuted] (3.2,2.6) -- (3.2,0.55);
        \node[text=cOurs, anchor=north, font=\tiny] at (1.9,0.25) {直せる部分（射影）};
        \node[text=mMuted, anchor=west, font=\tiny, align=left] at (3.3,1.6) {直せない\\部分};
        \foreach \p in {(0.7,0.25),(1.6,0.7),(2.4,0.1),(3.4,0.85)} \fill[cOurs] \p circle (1.5pt);
      \end{tikzpicture}\par
      {\raggedright\tiny 模式図．青の点はラベルの付いた発話\par}
    \end{column}
    \begin{column}{0.54\textwidth}
      \begin{itemize}\footnotesize\setlength{\itemsep}{10pt}
        \item $E_0 = W^{*} - W_0$：出発点の向きの誤差（$W^{*}$ はプールの全ラベルで学んだ生徒）
        \item \textbf{命題 1}：生徒の解の $W - W_0$ は，ラベルの付いた $4K$ 本の点が張る空間にある
        \item 直せるのは $E_0$ の高々 $4K$ 次元への射影だけ．どの部分かはラベルを付ける発話で決まる
      \end{itemize}
    \end{column}
  \end{columns}
  \hosoku{s:p2:span}{命題 1 の証明}
\end{frame}

\begin{frame}{ラベルを見ずに選んでも，向きの誤差を捉える発話は選べない（命題 2）}
  \rl{考え方と結果}
  \vspace{-0.2em}
  \begin{columns}[T,onlytextwidth]
    \begin{column}{0.55\textwidth}
      \centering
      \begin{tikzpicture}
        \begin{axis}[width=5.6cm, height=5.3cm, xmin=0.25, xmax=2.5, ymin=-0.6, ymax=7.6,
            ytick={0,...,7}, yticklabels={教える選び方（ラベルを使う）,白色化したノルム最大,白色化した k-center,core-set,TypiClust,予測で herding,余裕が最小,でたらめ（プールの順）},
            yticklabel style={font=\tiny}, xlabel={$E_0$ を捉える割合（でたらめ ＝ 1）},
            xmajorgrids, grid style={mMuted!25},
            legend style={font=\tiny, draw=none, at={(0.5,1.02)}, anchor=south, legend columns=3, /tikz/every even column/.append style={column sep=0.8em}}]
          \draw[black, densely dashed] (axis cs:1,-0.6) -- (axis cs:1,7.6);
          \draw[mMuted] (axis cs:0.25,0.5) -- (axis cs:2.5,0.5);
          \addplot[only marks, mark=*, mark size=1.6pt, kA, error bars/.cd, x dir=both, x explicit, error mark=none, error bar style={kA, line width=0.6pt}]
            table[x=a, y expr=\thisrow{y}+0.22, x error plus expr=\thisrow{a_hi}-\thisrow{a}, x error minus expr=\thisrow{a}-\thisrow{a_lo}] {capture.csv}; \addlegendentry{A}
          \addplot[only marks, mark=square*, mark size=1.5pt, kV, error bars/.cd, x dir=both, x explicit, error mark=none, error bar style={kV, line width=0.6pt}]
            table[x=v, y=y, x error plus expr=\thisrow{v_hi}-\thisrow{v}, x error minus expr=\thisrow{v}-\thisrow{v_lo}] {capture.csv}; \addlegendentry{V}
          \addplot[only marks, mark=triangle*, mark size=1.9pt, kT, error bars/.cd, x dir=both, x explicit, error mark=none, error bar style={kT, line width=0.6pt}]
            table[x=t, y expr=\thisrow{y}-0.22, x error plus expr=\thisrow{t_hi}-\thisrow{t}, x error minus expr=\thisrow{t}-\thisrow{t_lo}] {capture.csv}; \addlegendentry{T}
        \end{axis}
      \end{tikzpicture}\par
      {\raggedright\tiny $K=5$（20 本）．印はターゲットの平均どうしの比，線はターゲットごとの比の範囲．A 7，V 6，T 4 ターゲット，反復 10 回\par}
    \end{column}
    \begin{column}{0.43\textwidth}
      \begin{itemize}\scriptsize\setlength{\itemsep}{8pt}
        \item \textbf{命題 2}：正規近似のもとでは，ラベルを見ずに選んだどの $4K$ 本も，$E_0$ を捉える割合の期待値は同じ
        \item 実測：ラベル無しの 7 通りはでたらめの \NUM{p2-e58-ratio-unlab-lo}〜\NUM{p2-e58-ratio-unlab-hi} 倍．候補のラベルを使っても \NUM{p2-e58-ratio-lll-lo}〜\NUM{p2-e58-ratio-lll-hi} 倍
        \item 5 つの値の生徒では，どのラベル無しの選び方もでたらめを 1 \pt 上回らない
      \end{itemize}
    \end{column}
  \end{columns}
  \hosoku{s:p2:capture}{命題 2 の説明と選び方の定義}
\end{frame}

\begin{frame}{弱いラベル：先生の予測を，感情ごとの合計がそろうように割り直して全体に付ける}
  \rl{手法：弱いラベル}
  \vspace{-0.1em}
  \centering\scriptsize
  \setlength{\tabcolsep}{3.5pt}
  \begin{tabular}{@{}lcccc@{}}
    \multicolumn{5}{c}{\textbf{先生の予測 $Q$（例）}} \\
    \toprule
    & 中立 & 喜び & 悲しみ & 怒り \\
    \midrule
    発話 1 & 0.70 & 0.10 & 0.10 & 0.10 \\
    発話 2 & 0.60 & 0.20 & 0.10 & 0.10 \\
    発話 3 & 0.50 & 0.10 & 0.30 & 0.10 \\
    発話 4 & 0.10 & 0.20 & 0.10 & 0.60 \\
    \midrule
    列の合計 & \bad{1.90} & 0.60 & 0.60 & 0.90 \\
    \bottomrule
  \end{tabular}
  \quad\raisebox{-1.2em}{\begin{tabular}{c}$\Longrightarrow$\\[-0.2em]{\tiny 列ごとに}\\[-0.4em]{\tiny 合計で割る}\end{tabular}}\quad
  \begin{tabular}{@{}lcccc@{}}
    \multicolumn{5}{c}{\textbf{割り直した弱いラベル}} \\
    \toprule
    & 中立 & 喜び & 悲しみ & 怒り \\
    \midrule
    発話 1 & 0.37 & 0.17 & 0.17 & 0.11 \\
    発話 2 & 0.32 & 0.33 & 0.17 & 0.11 \\
    発話 3 & 0.26 & 0.17 & 0.50 & 0.11 \\
    発話 4 & 0.05 & 0.33 & 0.17 & 0.67 \\
    \midrule
    列の合計 & \good{1} & \good{1} & \good{1} & \good{1} \\
    \bottomrule
  \end{tabular}\par
  {\raggedright\tiny 説明用の数（実験の値ではない）．右の表は小数 2 桁に丸めた．このあと全体を $m/4$ 倍して合計を $m$ にする\par}
  \vspace{0.3em}\raggedright
  \begin{columns}[T,onlytextwidth]
    \begin{column}{0.49\textwidth}
      \begin{itemize}\scriptsize\setlength{\itemsep}{3pt}
        \item 先生は，ラベル無しの発話すべてに 4 感情の確率を出す
        \item 例の先生は中立を多く予測する（合計 1.90）．割り直すと，どの感情も合計が同じになる
      \end{itemize}
    \end{column}
    \begin{column}{0.49\textwidth}
      \begin{itemize}\scriptsize\setlength{\itemsep}{3pt}
        \item 生徒は $K$ 本と弱いラベルの両方で学ぶ．先生と $m$ は $K$ ごとに選ぶ
        \item 出発点自身の予測を割り直すだけでも先生の役をする
      \end{itemize}
    \end{column}
  \end{columns}
  \hosoku{s:p2:weak}{弱いラベルの作り方の詳細}
\end{frame}

\begin{frame}{同じモダリティの先生でも生徒は上がるが，その先生の $J$ は小さい}
  \rl{結果：見込みの蒸留}
  \vspace{0.1em}
  \begin{columns}[T,onlytextwidth]
    \begin{column}{0.44\textwidth}
      \centering
      \begin{tikzpicture}[font=\scriptsize]
        \node[font=\scriptsize\bfseries] at (2.0,3.0) {$J$ の定義（プールの発話で数える）};
        \node[font=\tiny] at (1.3,2.4) {出発点が正しい}; \node[font=\tiny] at (3.7,2.4) {出発点が誤る};
        \node[font=\tiny, anchor=east, align=right] at (0.05,1.6) {先生が\\正しい}; \node[font=\tiny, anchor=east, align=right] at (0.05,0.5) {先生が\\誤る};
        \draw[mMuted] (0.1,-0.05) rectangle (4.9,2.15); \draw[mMuted] (2.5,-0.05) -- (2.5,2.15); \draw[mMuted] (0.1,1.05) -- (4.9,1.05);
        \fill[mGood!25] (2.52,1.07) rectangle (4.88,2.13); \node[align=center, font=\tiny] at (3.7,1.6) {\good{\textbf{＋}}\quad 先生だけが当てる};
        \fill[mBad!20] (0.12,-0.03) rectangle (2.48,1.03); \node[align=center, font=\tiny] at (1.3,0.5) {\bad{\textbf{−}}\quad 出発点だけが当てる};
      \end{tikzpicture}\par
      \vspace{0.3em}
      {\raggedright\tiny $J$ ＝（先生だけが当てる割合）−（出発点だけが当てる割合）を感情ごとに平均したもの ＝ プールでの先生の UAR − 出発点の UAR\par}
    \end{column}
    \begin{column}{0.54\textwidth}
      \begin{itemize}\footnotesize\setlength{\itemsep}{10pt}
        \item 先生の情報から来る利得は $J$ を超えないと予想する（定理ではない）
        \item 同じモダリティの見込みの $J$ は，A で \NUM{p2-e52-q-minus-p-a} \pt しかない
        \item その弱いラベルで A は $K=0$ に \NUM{p2-e59k-a-k0} \pt 上がる．うち \NUM{p2-e59k-ctrl-a-k0} \pt は釣り合い直しだけの分
      \end{itemize}
    \end{column}
  \end{columns}
  \hosoku{s:p2:e59}{見込みの蒸留の詳細}
\end{frame}

\begin{frame}{モダリティをまたぐ先生：ソースだけで学んだ融合の出発点が生徒に教える}
  \rl{手法：別のモダリティの先生}
  \vspace{0.1em}
  \centering
  \begin{tikzpicture}[font=\scriptsize]
    \node[anchor=west, font=\scriptsize\bfseries] at (-0.2,2.55) {単独の生徒（例：AV の音声の生徒）};
    \node[draw=cExt, fill=white, rounded corners=2pt, align=center, text width=2.9cm] (f1) at (1.4,1.75) {融合の出発点（AV）\\ソースだけで学び，学び直さない};
    \node[draw=cExt, fill=cExt!8, rounded corners=2pt, align=center, text width=2.0cm] (w1) at (4.9,1.75) {弱いラベル\\（プール全体）};
    \node[draw=cOurs, fill=white, rounded corners=2pt, align=center, text width=2.0cm] (a1) at (7.8,1.75) {音声の生徒};
    \node[testbox, text width=2.3cm] (t1) at (10.8,1.75) {テストでは音声だけ};
    \draw[arr, cExt] (f1) -- (w1); \draw[arr, cExt] (w1) -- (a1); \draw[tarr] (a1) -- (t1);
    \node[anchor=west, font=\scriptsize\bfseries] at (-0.2,0.75) {融合の生徒（例：AV の生徒）};
    \node[draw=cExt, fill=white, rounded corners=2pt, align=center, text width=2.9cm] (f2) at (1.4,-0.15) {融合の出発点（AV）\\ソースだけで学ぶ};
    \node[draw=cOurs, fill=white, rounded corners=2pt, align=center, text width=2.4cm] (a2) at (5.0,0.3) {音声の生徒（プールで学ぶ）};
    \node[draw=cOurs, fill=white, rounded corners=2pt, align=center, text width=2.4cm] (v2) at (5.0,-0.6) {顔の生徒（プールで学ぶ）};
    \node[draw=cExt, fill=cExt!8, rounded corners=2pt, align=center, text width=2.0cm] (p2) at (8.2,-0.15) {2 つの予測の積\\を弱いラベルに};
    \node[draw=cOurs, fill=white, rounded corners=2pt, align=center, text width=2.0cm] (s2) at (11.0,-0.15) {AV の生徒};
    \draw[arr, cExt] (f2.east) -- node[above, font=\tiny, text=cExt, pos=0.45] {教える} (a2.west);
    \draw[arr, cExt] (f2.east) -- (v2.west);
    \draw[arr, cOurs] (a2.east) -- (p2.west); \draw[arr, cOurs] (v2.east) -- (p2.west);
    \draw[arr, cExt] (p2) -- (s2);
  \end{tikzpicture}\par
  \vspace{0.3em}\raggedright
  \begin{itemize}\scriptsize\setlength{\itemsep}{4pt}
        \item 融合の出発点は，単独のモダリティとは違う向きを知っている（特権情報と同じ考え方）
        \item 下調べ（AV）の $J$：融合の出発点は，音声の出発点に対して \NUM{p2-e60pre-juar-av-a-fused} \pt
        \item 基本版の先生の候補：融合の出発点，単独の出発点の予測の積，2 つの一致で重みづけた融合の出発点（どれもソースだけで学ぶ）
      \end{itemize}
  \hosoku{s:p2:pre}{下調べの $J$ の全表}
\end{frame}

\begin{frame}{$K\le5$ では，モダリティをまたぐ弱いラベルがすべての生徒を出発点より上げる}
  \rl{結果：モダリティをまたぐ蒸留}
  \vspace{-0.2em}
  \begin{columns}[T,onlytextwidth]
    \begin{column}{0.48\textwidth}
      \centering
      \begin{tikzpicture}
        \begin{axis}[width=6.8cm, height=4.7cm, xmin=-0.2, xmax=4.2, ymin=-1, ymax=4.5,
            xtick={0,...,4}, xticklabels={0,1,5,10,25}, xlabel={感情あたりのラベル $K$},
            ylabel={出発点との差（\pt）}, ymajorgrids, grid style={mMuted!25},
            legend style={font=\tiny, draw=none, at={(0.5,1.02)}, anchor=south, legend columns=2, /tikz/every even column/.append style={column sep=0.5em}}]
          \draw[black, line width=0.6pt] (axis cs:-0.2,0) -- (axis cs:4.2,0);
          \addplot[mMuted, mark=o, line width=1.1pt] table[x=x, y=e59] {ladder_gains.csv}; \addlegendentry{同じモダリティの先生（7）}
          \addplot[cExt!55, mark=diamond*, line width=1.1pt] table[x=x, y=e60av] {ladder_gains.csv}; \addlegendentry{AV の基本版（6）}
          \addplot[cExt, mark=square*, line width=1.1pt] table[x=x, y=e60avt] {ladder_gains.csv}; \addlegendentry{AVT の基本版（4）}
          \addplot[cExt, mark=*, line width=1.4pt, densely dashed] table[x=x, y=oursx] {ladder_gains.csv}; \addlegendentry{拡張版，AVT（4）}
        \end{axis}
      \end{tikzpicture}\par
      {\raggedright\tiny 音声の生徒．括弧はターゲットの数．線ごとにターゲットと出発点が違うので，高さでなく出発点との差を比べる．反復 10 回\par}
    \end{column}
    \begin{column}{0.5\textwidth}
      \begin{itemize}\footnotesize\setlength{\itemsep}{10pt}
        \item 単独の生徒 \NUM{p2-e60-m1-n} 個は，どれも $K=0$ と $K=5$ で上がる（$K=0$ の AV の A は \NUM{p2-e60-av-a-k0} \pt）
        \item 融合の生徒は $K=0$ で \NUM{p2-e60-f-k0-lo}〜\NUM{p2-e60-f-k0-hi} \pt 上がる
        \item $K\ge10$ では \NUM{p2-e60-m3-students-n} の生徒のうち \NUM{p2-e60-m3-loss-n} が出発点を下回る
      \end{itemize}
    \end{column}
  \end{columns}
  \hosoku{s:p2:e60}{全生徒・全 $K$ の値}
\end{frame}

\begin{frame}{拡張版は $K$ 本で学び直した融合の生徒を先生にして，$K\ge10$ の損を消す}
  \rl{手法：基本版と拡張版}
  \vspace{0.3em}
  \centering\scriptsize
  \renewcommand{\arraystretch}{1.35}
  \begin{tabularx}{\textwidth}{@{}l X X@{}}
    \toprule
    & \textbf{基本版} & \textbf{拡張版} \\
    \midrule
    単独の生徒の先生の候補 & 融合の出発点，単独の出発点の予測の積，2 つの一致で重みづけた融合の出発点（どれもソースだけで学ぶ） & $K$ 本で学び直した融合の生徒（$K=0$ では融合の出発点），それが教えた単独の生徒の予測の積 \\
    融合の生徒の先生の候補 & 単独の出発点の予測の積，融合の出発点が教えた単独の生徒の予測の積（弱いラベルにするか，そのまま予測に使う） & $K$ 本で学び直した融合の生徒が教えた単独の生徒の予測の積 \\
    弱いラベルの重みづけ & 感情ごとの合計をそろえる & そろえる，そろえない，しきい値を超えた発話だけを硬いラベルにする \\
    $K$ ごとに選ぶもの & 先生，$m$ & 先生，$m$，重みづけ \\
    \bottomrule
  \end{tabularx}\par
  \vspace{0.6em}\raggedright
  \begin{itemize}\scriptsize\setlength{\itemsep}{4pt}
        \item 拡張版は，AVT の音声の生徒と AV の融合の生徒で $K\ge10$ の損を消した
        \item どちらの版でも，先生になる単独の生徒はプールのラベル無しの発話で学び，$K\ge1$ では $K$ 本でも学ぶ
      \end{itemize}
  \hosoku{s:p2:ext}{拡張版の候補}
\end{frame}

\begin{frame}{第 II 部のまとめ：向きの情報は，別のモダリティの弱いラベルから得る}
  \rl{まとめ}
  \vspace{0.3em}
  \begin{itemize}\footnotesize\setlength{\itemsep}{9pt}
        \item 出発点は，中心化，全ソースからの仮説転移，5 つの値の選択の順に UAR が上がる
        \item 少ないラベルが直せる向きは $4K$ 本の張る空間の中だけで，ラベルを見ない選び方ではそこを向きの誤差に合わせられない
        \item 同じモダリティの先生は優位 $J$ が小さい
        \item ソースだけで学んだ別のモダリティの先生（融合の出発点など）を使うと，$K\le5$ ですべての生徒が出発点を上回る
      \end{itemize}
\end{frame}
